Tras un arranque exitoso, los servicios quedan accesibles en los puntos
recogidos en la \cref{table:install-services}. % TODO cuadro en vez de tabla

\begin{table}[Servicios expuestos tras el arranque]%
            {table:install-services}%
            {Servicios expuestos tras el arranque y credenciales por
             defecto.}
  \small
  \renewcommand{\arraystretch}{1.3}
  \begin{tabular}{l l l}
    \hline
    \textbf{Servicio} & \textbf{\acs{uri}} & \textbf{Credenciales por defecto} \\
    \hline
    \acs{api} \acs{rest}      & \texttt{http://localhost:8000/api/v1/}                  & (las del superusuario) \\
    Swagger UI                & \texttt{http://localhost:8000/api/schema/swagger-ui/}   & --- \\
    MinIO Console             & \texttt{http://localhost:9001}                          & \texttt{minioadmin / minioadmin} \\
    Mailpit                   & \texttt{http://localhost:8025}                          & --- \\
    \hline
  \end{tabular}
\end{table}

Los servicios de monitorización del perfil \texttt{monitoring}
(pgAdmin, Redis Commander, Flower) se levantan opcionalmente con
\texttt{docker compose -{}-profile monitoring up -d}.
